\section{Delivering Ahead of the User: Adaptation, Prediction and Caching}\label{sec:anticipate}

An immersive client displays only a fraction of what it could receive, and it must request that fraction before the user looks there. This section asks how a system decides what to fetch, at what quality and from where, ahead of the user's head and body. Three problem cells answer the question. Rate adaptation and adaptive streaming is the largest problem cell in the atlas, with \nxrNProbRateAdaptationStreaming{} papers (\nxrShProbRateAdaptationStreaming{} of the corpus), but its share fell from \nxrShProbRateAdaptationStreamingEarly{} in 2017--2019 to \nxrShProbRateAdaptationStreamingLate{} in 2022--2024: it grew by \nxrCagrProbRateAdaptationStreaming{} a year while the corpus grew by \nxrCagrTotal{}. Viewport and motion prediction (\nxrNProbViewportMotionPrediction{} papers, \nxrCagrProbViewportMotionPrediction{} a year) and caching and prefetching (\nxrNProbCachingPrefetching{} papers, \nxrCagrProbCachingPrefetching{} a year) are smaller and growing faster. The decline mirrors that of 360-degree video, the content all three were built for, whose share fell from \nxrShContVideoThreeSixtyEarly{} to \nxrShContVideoThreeSixtyLate{} over the same periods. Evidence thins along the chain: \nxrRealRateAdaptationStreaming{} of rate-adaptation papers use a real testbed and \nxrUsersRateAdaptationStreaming{} involve users; prediction is judged mainly on recorded traces (\nxrTraceViewportMotionPrediction{}); caching is judged mainly in simulation (\nxrSimCachingPrefetching{}) and never with users. Representations and codecs are treated in \cref{sec:encode}, and remote rendering in \cref{sec:compute}.

\subsection{Viewport-adaptive streaming: spending bits where the user looks}

Tiled delivery of 360-degree video saves bandwidth, and the saving shrinks as the test becomes realistic (\cref{tab:anticipate-savings}). With a scripted viewport and an all-high-quality baseline, hexaface tiling saved up to 72\% \cite{W2524582890}; HEVC tiles saved 65\% for a static viewport and over 40\% with recorded head motion \cite{W2623784065}; Flare, run on phones over commercial LTE, saved 26--35\% \cite{W2896840651}. The saving is paid for twice. Tiling costs coding efficiency, a median 7.3--10.4\% for 1-s tiles \cite{W2896840651} and up to 37.3\% more bitrate for an 8$\times$5 grid at low resolution \cite{W2623784065}. Decoding many tiles costs the client: Flare raised phone CPU load from 18\% to 58\% \cite{W2896840651}, and Rubiks layered its tiles to fit smartphone decoders \cite{W2826107421}. Prediction error is the main limiter, so controllers hedge. 360ProbDASH fetched tiles in proportion to their viewing probability and raised viewport PSNR by about 39\% \cite{W2766324256}, and segments of at most 2--3~s bound the cost of a wrong guess \cite{W2527384707}.

How the controller should be built is contested. Deep reinforcement learning, used by about a quarter of the cell, reports 20--30\% gains over earlier heuristics \cite{W2920378308,W3159074940}; BOLA360, a buffer-based controller with a near-optimality proof and no training, reports 6--110\% \cite{W4394890851}. The two families have never been compared on shared traces, and learned policies degrade on unseen viewing patterns and networks \cite{W4403826584,W4400409717}. Both rest on an assumed objective: a hand-weighted sum of viewport quality, rebuffering and quality switches. No abstract in the cell names an ITU-T QoE model, although tile extensions of P.1203 and P.1204.3 exist \cite{W3200608537,W3168206114}. The quality term is itself disputed: one study favours PSNR \cite{W2741833902}, another VMAF \cite{W3115841380}, and a database of 378 viewport videos found that mainstream models fail for viewport-adaptive streams \cite{W4402916896}. Even delay sensitivity is unsettled. An OMAF analysis found end-to-end delay decisive for tile adaptation \cite{W2914790472}, whereas a subjective test found that 47~ms of delay barely changed ratings \cite{W2949677285}.

Most of this evidence comes from one on-demand client on an emulated link. Live streaming removes the pre-encoded ladder and brings the uplink into the loop. Live360 sets camera upload rates by viewer interest and reports twice the average QoE of its baselines, in its own metric \cite{W4317796966}, while YouTube Live delivered 360-degree streams with 4.2--44.2~s camera-to-display delay depending on mode \cite{W4361857501}. Multicast exploits the similarity of viewers' motion and saved over 50\% \cite{W2725312916}, but adaptation for several headsets behind one bottleneck has reached a testbed in only three papers and users in none \cite{W4377000454,W2890617053}. Wireless access is the largest blind spot. Only two abstracts in the cell use commercial cellular traces \cite{W3021314748,W2810704880}, although mmWave variability and handover already require new mechanisms such as band switching \cite{W3195004621}.

\begin{table}[t]
\centering\small
\caption{Bandwidth saved by viewport- and visibility-adaptive delivery, ordered within each content type from idealised to realistic evaluation. The saving falls as the viewport and the network become real.}
\label{tab:anticipate-savings}
\begin{tabularx}{\linewidth}{@{}lXrl@{}}
\toprule
Content & Condition & Saving & Source\\
\midrule
360-degree & Perfect knowledge of future motion (bound) & 80\% & \cite{W2585813514}\\
360-degree & Scripted viewport, against all tiles at high quality & up to 72\% & \cite{W2524582890}\\
360-degree & HEVC tiles, static viewport / recorded head motion & 65\% / $>$40\% & \cite{W2623784065}\\
360-degree & Predicted motion with redundancy, $<$0.1\% failure & 45\% & \cite{W2585813514}\\
360-degree & Phones on commercial LTE, 4 users, 3 videos & 26--35\% & \cite{W2896840651}\\
Volumetric & Gaze-driven delivery on HoloLens~2 & up to 67\% & \cite{W4399323898}\\
Volumetric & Visibility-aware, commercial 5G, headset / phone users & 40\% / 10\% & \cite{W3021314748}\\
\bottomrule
\end{tabularx}
\end{table}

\subsection{Six degrees of freedom: adapting volumetric scenes}

Volumetric streaming carries the same idea into 6DoF, where the unit of adaptation becomes a 3D cell, an object or a level of detail, and visibility depends on position, distance and occlusion as well as direction. It is the fastest-growing part of rate adaptation, at about forty papers a year by 2025. ViVo combined the three visibility tests and saved 40\% of data on average, and up to 80\%, over commercial 5G at SSIM above 0.99 \cite{W3021314748}. The same content saved 40.0\% for headset users but only 9.98\% for phone users, so the gain depends on how people move. Related designs allocate rate across 3D tiles with a window-based buffer \cite{W2963389850}, solve a QoE-driven knapsack \cite{W4210824758}, keep octrees decodable in parallel on mobile devices \cite{W3088533413}, or refine quality progressively with neural compression \cite{W4392405518}.

No general allocation policy has emerged. PCC-DASH found that the best multi-object strategy depends on the scene and the camera path \cite{W2982262421}, and its subjective follow-up had to test whether objective and subjective rankings even agree \cite{W3036519337}. The field assumes that its navigation traces are representative, yet the newest public set has 15 participants \cite{W4408870156}, and only about one volumetric adaptation paper in thirty evaluates with users. Two-way volumetric communication breaks the buffer-based model altogether. Systems reuse low-latency DASH \cite{W3030878847} or WebRTC \cite{W4394890952}, a tiled user-adaptive point-cloud system needed 14~Mbps instead of about 300~Mbps in a 33-user study \cite{W4304091566}, and LiVo reached a mean opinion score of 4.1 \cite{W4416659451}. In these systems, adaptation is coupled loosely, if at all, to congestion control. Gaussian splatting has layered representations meant for adaptive streaming, but they are judged on visual fidelity rather than on behaviour under variable throughput \cite{W4402704690}.

\subsection{Prediction: how far ahead, and against which baseline}

The prediction horizon couples the two halves of the controller. A deep buffer absorbs throughput dips but forces the system to predict the viewport seconds ahead; a shallow buffer needs only short predictions but stalls when throughput drops. Accuracy falls steeply with horizon: Flare's best predictor reached median accuracies of 90.5\%, 72.9\%, 58.2\% and 35.2\% at 0.2, 0.5, 1 and 3~s \cite{W2896840651}, and a cross-dataset benchmark found no algorithm accurate much beyond 300~ms \cite{W3089527252}. For on-demand content, earlier viewers are the strongest long-horizon signal. Cross-user prediction raised accuracy by 20.2 percentage points \cite{W2896449077}, and adding other users' viewports raised the 1--10~s hit rate from 0.7283 to 0.7791 \cite{W2963428739}. Live content has no such history and relies on online learning or object semantics \cite{W4243538177,W3141438256}. Saliency and gaze help in ablations, more at long horizons \cite{W4387148495,W4387968308}.

Whether deep models beat simple predictors is the central controversy. On the Tsinghua dataset, persistence (keep the current viewport) reached a hit rate of 0.8355, against 0.8402 for a sequence-to-sequence model and 0.8464 for a single LSTM \cite{W2963428739}; for 6DoF pose, simple statistical models performed surprisingly well \cite{W4379805466}. Learned models report large gains mostly over other learned models. A fusion transformer reached an average precision of 0.954 at 1~s against 0.901 for the previous best, but on a random split of sliding windows that placed the same videos and users in training and test \cite{W4387148495}. Its accuracy also had a price, 74~ms per batch on a desktop GPU against 21~ms for a scanpath-only transformer \cite{W4387148495,W4251137679}. Fewer than one prediction paper in ten names a persistence, linear or Kalman baseline, and no evaluation protocol with disjoint videos and users is accepted.

Two further assumptions go untested. Controllers need calibrated uncertainty to size tile margins, yet predictors output point estimates. Predicting the predictor's own deviation saved 45\% of bandwidth at under 0.1\% failure \cite{W2585813514}, and some models output several possible futures \cite{W4280601181}, but no paper the atlas read reports calibration. Second, accuracy is assumed to translate into experience, yet closed-loop tests with real systems fell from 9 of 74 prediction papers in 2016--2022 to 1 of 54 in 2023--2026. Wireless uses of prediction inherit the problem. A federated predictor for mmWave association reduced breaks in presence by up to 16\% and 26\% against centralised baselines, using traces that paired hourly location samples with unrelated head orientations \cite{W2975296300}, whereas a 60~GHz measurement found 3D motion too unpredictable for beam steering \cite{W2791171246}. Pose prediction that hides motion-to-photon (MTP) latency in remote rendering is discussed in \cref{sec:compute}.

\subsection{Caching and prefetching: placing content before it is requested}

Viewport-aware caching beats popularity-only replacement, a result replicated by independent groups but always against weak baselines. Placing tiles by cross-user viewport similarity raised the edge hit ratio by at least 40\% over LFU and 17\% over LRU \cite{W2896595273}. OpCASH covered more than 95\% of the viewport from cache after 24 views and fetched 85\% less data from the origin \cite{W4288370107}; VRCT rebuilt viewports from partly cached tiles, raising hit ratio by 29\% and cutting backhaul by 44\% \cite{W4377008005}; CUBIST combined reactive and proactive caching for an 87\% hit ratio \cite{W3213361583}. For live 360-degree video, predictive edge caching cut backhaul by up to 76\% \cite{W4390136482}. Prefetching extends to AR and volumetric content: world-scale AR retrieval missed 20\% fewer deadlines with ten times fewer wasted bytes \cite{W4379932440}, and PACE raised modelled QoE by up to 54.8\% \cite{W4415047817}. Joint designs that cache and also transcode, render or synthesise views at the edge report lower latency in simulation \cite{W2979447669,W4210795263}, a result that holds almost by construction, because the joint problem contains the decoupled one.

Three questions remain open. The first is granularity: fixed tiles, variable tiles that rarely recur \cite{W4288370107}, or similarity caching over continuous orientations \cite{W3156984197}, never compared head to head. The second is the interaction with adaptation: cache hits distort the throughput estimates that DASH clients rely on \cite{W3122351391}, yet most caching papers assume an ideal client. The third is whether multi-agent and graph learners add anything beyond good viewport prediction, which cannot be settled while XRFlux remains the only shared benchmark \cite{W4409763640}. Wireless caching sharpens the gap between theory and practice. Multi-antenna coded caching promises 100--400\% more parallel streams, but its authors list subpacketisation and synchronised requests as obstacles \cite{W4381251523}, and a learned UAV small-cell cache was tested with 20 simulated users \cite{W2962896865}. The common assumption is that hit ratio stands for experience. No caching scheme in the atlas has been evaluated with users, and fewer than one caching paper in twenty addresses 6DoF content, which is reusable across viewpoints and users.

\begin{openproblems}
\item \textbf{Prediction and adaptation are not validated in closed loop with users.} Real-system tests of predictors fell to 1 of 54 papers after 2022, and ABR objectives are hand-weighted. A tiled player that runs persistence, Kalman and learned predictors under ABR policies over real bandwidth, with a within-subjects rating study, would show whether accuracy rankings survive.
\item \textbf{No dataset joins head motion, network conditions and ratings from the same session.} Four dataset papers pair head data with network traces and four with ratings; replay removes the viewer's reaction to quality. Live sessions over Wi-Fi and commercial 5G, logging pose, packets and continuous ratings, would close this gap.
\item \textbf{Adaptation over real wireless access is under-evaluated.} The cell has half the expected 5G papers and two abstracts with commercial cellular traces. Paired head-motion and radio traces from walking and seated users, replayed in emulation and run live, would test whether PHY and MAC hints reduce stalls.
\item \textbf{6DoF adaptation and caching lack user-validated QoE and realistic navigation traces.} Volumetric streaming grows fastest, yet trace sets have about 15 participants and volumetric caching is almost unstudied. Larger 6DoF traces, paired comparisons of rendered views and a level-of-detail edge cache shared by co-located users are needed.
\item \textbf{Caching gains have never been shown to users or outside simulation.} No caching paper has users, and none of the coded, D2D or UAV caching papers uses traces or a testbed. An unmodified HTTP edge cache serving tiled DASH, comparing LRU, viewport-aware and Belady policies with viewers, would quantify perceptible benefit.
\item \textbf{Motion-driven wireless decisions use mismatched position and orientation traces.} Seven of nine such papers are simulations, and the one measurement found motion too unpredictable for 60~GHz beam steering. Room-scale multi-user sessions recording 6DoF tracking with synchronised link measurements would show how far ahead blockage is predictable.
\end{openproblems}
